% TOP_PROBLEM: {"id": 2305034, "title": "Research Problems in Function Theory — Problem 5.34", "category_id": 8, "category": {"id": 8, "name": "analysis", "display_name": "Analysis", "description": "Limits, continuity, calculus, and function theory.", "slug": "analysis", "order_index": 8, "created_at": "2026-07-31T15:26:25.671Z"}, "status": "open"}
% FIRST_POSTED: null
% ENTRY_KIND: statement_only
% Imported from: batch14.tex; UnsolvedMath ID 2305034
% Compile twice: lualatex 2305034.tex
\documentclass[11pt]{article}
\usepackage{fontspec}
\setmainfont{Latin Modern Roman}
\usepackage{amsmath,amssymb,mathtools,unicode-math}
\setmathfont{Latin Modern Math}
\usepackage{xurl,hyperref}
\hypersetup{hidelinks,pdftitle={UnsolvedMath Problem 2305034}}
\newcommand{\sref}[2]{\hyperlink{source:#1:#2}{[S#2]}}
\newcommand{\cataloguescope}[1]{\paragraph{Mathematical question.}#1}
\begin{document}
% UnsolvedMath ID 2305034; source catalogue code AMR-022-5034
\setcounter{section}{2305033}
\section{Research Problems in Function Theory — Problem 5.34}\label{2305034}
{\small\sffamily UnsolvedMath ID 2305034 · Analysis}
\subsection{Definitions and mathematical statement}

\paragraph{Shared notation.}$\mathbb N=\{1,2,\ldots\}$, and $\mathbb Z,\mathbb Q,\mathbb R,\mathbb C$ denote the integers, rationals, reals and complex numbers. Any other conventions are specified locally.


Write $\mathbb D=\{z\in\mathbb C:|z|<1\}$ and $\mathbb T=\partial\mathbb D$. All Taylor coefficients below are taken at $0$. For $\zeta\in\mathbb T$, a Stolz angle means a nondegenerate Euclidean angular sector with vertex $\zeta$, whose two rays enter $\mathbb D$ nontangentially, restricted to a sufficiently small neighborhood of $\zeta$. A finite angular limit is a common finite limit as $z\to\zeta$ within every such sector.A Bloch function is holomorphic $h$ with $\sup_{z\in\mathbb D}(1-|z|^2)|h'(z)|<\infty$. Let $E_h\subset\mathbb T$ be the set where $h$ has an angular limit in $\widehat{\mathbb C}$, allowing $\infty$. A Hausdorff gauge is a nondecreasing function $\psi:[0,\infty)\to[0,\infty)$ with $\psi(0)=0$, $\psi(t)>0$ for $t>0$, and $\psi(t)\to0$ as $t\downarrow0$. Define
\[\mathcal H^\psi(E)=\lim_{\varepsilon\downarrow0}\inf\left\{\sum_j\psi(\operatorname{diam}U_j):E\subseteq\bigcup_j U_j,\ \operatorname{diam}U_j\le\varepsilon\right\}.
\]
For the logarithmic gauge take $\psi_{\log}(t)=1/\log(1/t)$ for sufficiently small $t>0$; its values away from $0$ do not affect the resulting Hausdorff measure.
\paragraph{Statement recorded in the supplied source.}
Is there a single Hausdorff gauge $\psi$, independent of the Bloch function, such that $\mathcal H^\psi(E_h)>0$ for every Bloch function $h$? In particular, does this hold for the logarithmic gauge $\psi_{\log}$? \sref{2305034}{2}
\subsection{Short English statement}
Can angular-limit points of every Bloch function be guaranteed positive measure for one universal Hausdorff gauge, particularly a logarithmic gauge?
\subsection{Sources}
\begin{description}
\item[\hypertarget{source:2305034:1}{S1}] ULAM.AI and the UnsolvedMath Contributors, UnsolvedMath: A Curated Collection of Open Mathematics Problems, record 2305034 (AMR-022-5034). CC BY 4.0. \url{https://huggingface.co/datasets/ulamai/UnsolvedMath}.
\item[\hypertarget{source:2305034:2}{S2}] W. K. Hayman and E. F. Lingham, Research Problems in Function Theory (2018), arXiv:1809.07200, Chapter 5, Problem 5.34, its update and the preceding shared notation. \url{https://arxiv.org/abs/1809.07200}.
\end{description}
\label{end:2305034}
\end{document}
